\documentclass[12pt]{article}
\usepackage[utf8]{inputenc}
\usepackage[T1]{fontenc}
\usepackage{hyperref}
\title{Yawarawarga Language Resources for K-12 Curriculum Development}
\author{Compiled by Claude Sonnet 5.5 (claude-sonnet-5-5) for Andr\'as Kornai}
\date{}
\begin{document}
\maketitle

\section{Introduction}
Yawarawarga (ISO 639-3: yww; also spelled Yawarrawarrka) is an extinct Australian Aboriginal language of the Pama-Nyungan family, associated with far western Queensland and the adjoining part of South Australia towards Innamincka. Sources disagree on its status: Wikipedia reports that Dixon (2002) treats it as a dialect of Yandruwandha, while Bowern (2001) treats it as a closely related language. In this search I found very little material specifically on Yawarawarga, and no curriculum resources, so this dossier is short.

\section{Primary Educational Resources}
I found no primers, textbooks or school programmes for Yawarawarga.

\section{Linguistic and Reference Materials}
\begin{itemize}
\item Wikipedia, Yawarrawarrka (overview of the language and its classification): \url{https://en.wikipedia.org/wiki/Yawarrawarrka}
\item Wikipedia, Yandruwandha language, the closely related language whose materials are the likeliest basis for any revival work: \url{https://en.wikipedia.org/wiki/Yandruwandha_language}
\item State Library of Queensland, ``Word of the Week'' blog post, week six (2019). It gives the name's etymology: \emph{yawara} `word or language' and \emph{waka} `small'. \url{https://www.slq.qld.gov.au/blog/2019-international-year-indigenous-languages-word-week-week-six}
\end{itemize}

\section{Religious and Community Resources}
I found no scripture translations or community language organisations specific to Yawarawarga.

\section{Digital and Audio Resources}
I found no audio or digital resources for this language.

\section{Curriculum Development Context}
The language is extinct and thinly documented. Any school programme would depend on consultation with descendant communities and on archival wordlists, and would probably have to draw on documentation of Yandruwandha. Teachers should verify all material with the relevant Aboriginal community before use.
\end{document}